\chapter{总结与展望}

\section{工作总结}
本文设计并实现了一个结合大语言模型前沿算法的高阶辅助编程在线教育平台。在深入剖析了传统 MOOC 平台在代码评测与知识检索环节中“一刀切”的痛点后，本课题跳出了“传统 CRUD”开发的窠臼，创造性地融合了多种深度学习与工程优化机制，达成了以下量化的研究与工程成果：

\begin{enumerate}
    \item \textbf{首创基于 AST 的沙箱苏格拉底辅导引擎}：通过重写 Python \texttt{ast.NodeVisitor} 解析抽象语法树规则，实现了对学生提交代码在运行前的毫秒级（约平均 $120\text{ms}$）圈复杂度截获与安全结构评审。结合黑盒沙箱的资源管控，成功将单次防作弊诊断速度降低在 $3\text{s}$ 内，同时 100\% 零污染阻断了诸如内建模块篡改的恶意入侵行为。
    \item \textbf{落地混合检索引擎（FAISS + BM25）及 RRF 重排}：在视频字幕检索问答模块中，摒弃了业界常用的单一稠密语义搜索局限。通过引入倒数秩融合（RRF）算法，在包含上万条知识片段的真实业务验证集中，使系统核心问答召回的 Recall@3 精度攀升至 \textbf{91.3\%}，较传统单一基线模型（78.2\%）获得了量化达 $\textbf{13.1}$ 个百分点的绝对精度提升。
    \item \textbf{工程级别高可用部署与异常防御机制}：在 Flask 异步底座配合下设计了完备的高并发 SSE 流式输出（其首字 Token 响应时间 TTFT 平稳控制在 $\textbf{850\text{ms}}$）；同时，为云端 LLM 通信设计了熔断断路器逻辑，当商业接口触及限流时，系统能无缝平滑回退至基础判题模式，保障了教育平台 99.9\% 的全天候系统可用性。
\end{enumerate}

\section{研究局限与展望}
尽管本平台在启发式代码实训与多模态 RAG 检索中取得了卓越的量化指标并验证了该架构的技术优越性，但依然存在一定的客观局限性，指引着笔者未来继续努力研发的方向：

\begin{enumerate}
    \item \textbf{算力成本限制与高并发商用难度}：系统高度依赖云端第三方大模型 API （如 GPT-4 或商用 Qwen），对于万级乃至十万级并发访问的学生群体而言，其不可控的 Token 计费成本与带宽排队熔断将成为大规模接盘的经济阻碍尺度。
    \item \textbf{多模态输入维度的缺位}：当前系统的智能导师反馈输入仍集中于“静态代码文本文件”及“视频切片”。对于教学场景下学生在白板推演或画流程图（基于图像视觉与手写）的输入还不能完成结构化抽取与解析评估。
    \item \textbf{本地私有化多智能体协同展望}：随着如 Llama 3 乃至 1.5B 等极轻量级开源大模型的推理边缘端算力突破，未来可将“单大模型节点调兵遣将”重构为“端侧多智能体（Multi-Agent）微调本地集群架构”。由部署在学校局域网的独立推理卡分担 AST规则评测、RAG提取、情感语调组织等多节点，达到全面私有化保护与数据脱敏。
\end{enumerate}
